\section{Illustrative New Keynesian Welfare Calibration}\label{app:nk_welfare}

This appendix reports a model-based welfare exercise that
maps the estimated overreaction parameter into a
three-equation New Keynesian benchmark. The exercise is
illustrative: it quantifies the order of magnitude of
welfare costs from diagnostic expectations under standard
calibration, not a policy-invariant welfare mapping.

The model is
\begin{align}
  \pi_t &= \beta E_t^\theta[\pi_{t+1}] + \kappa x_t + u_t, \nonumber \\
  x_t &= E_t[x_{t+1}] - \sigma^{-1}
    \left(i_t - E_t^\theta[\pi_{t+1}]\right), \nonumber \\
  i_t &= \phi_\pi \pi_t + \phi_x x_t, \nonumber
\end{align}
where $E_t^\theta$ denotes diagnostic expectations with
overreaction parameter $\theta$. Calibration uses standard
values: $\beta = 0.99$, $\kappa = 0.05$, $\sigma = 2$,
$\phi_\pi = 1.5$, $\phi_x = 0.5$, and AR(1) cost-push
shocks with persistence $\rho_u = 0.8$.

The welfare loss function is
$\mathcal{L}(\theta) = \mathrm{Var}(\pi) +
\lambda\,\mathrm{Var}(x)$ with
$\lambda = \kappa/\varepsilon \approx 0.0083$.
Table~\ref{tab:nk_welfare} reports welfare losses for a
range of $\theta$ values. At $\theta = 0.55$, aligned
with the household panel estimate, the welfare loss is
approximately 38\% above the rational benchmark
($\theta = 0$). At $\theta = 1.0$, a value closer to
what the aggregate point estimate would imply, the loss
rises to 78\% above benchmark.

\begin{table}[htbp]
\centering
\caption{Welfare Costs of Diagnostic Inflation Expectations}
\label{tab:nk_welfare}
\begin{tabular}{ccccc}
\toprule
$\theta$ & $\mathrm{Var}(\pi)$ & $\mathrm{Var}(x)$ & Loss $\mathcal{L}(\theta)$ & Welfare Cost (\%) \\
\midrule
  0.00 & 45.5716 & 27.5680 & 45.8013 & 0.00 \\
  0.25 & 52.6375 & 28.5526 & 52.8754 & 15.45 \\
  0.50 & 60.9654 & 29.5838 & 61.2119 & 33.65 \\
  0.55$^{\dagger}$ & 62.7824 & 29.7956 & 63.0307 & 37.62 \\
  0.75 & 70.5552 & 30.6616 & 70.8107 & 54.60 \\
  1.00$^{\ddagger}$ & 81.4071 & 31.7859 & 81.6720 & 78.32 \\
  1.20 & 90.9972 & 32.7188 & 91.2698 & 99.27 \\
  1.50 & 106.8967 & 34.1741 & 107.1815 & 134.01 \\
\bottomrule
\end{tabular}
\begin{tablenotes}
\small
\item \textit{Notes:} Welfare loss $\mathcal{L}(\theta) = \mathrm{Var}(\pi) + \lambda \, \mathrm{Var}(x)$ with $\lambda = \kappa/\varepsilon = 0.05/6 \approx 0.0083$.
\item Calibration: $\beta=0.99$, $\kappa=0.05$, $\sigma=2$, $\phi_\pi=1.5$, $\phi_x=0.5$, $\rho_u=0.8$.
\item Welfare cost is the percentage increase in loss relative to rational benchmark ($\theta=0$).
\item $^{\dagger}$ Household survey estimate ($\hat{\beta}_{\mathrm{rev}} \approx -0.55$). $^{\ddagger}$ Aggregate implied ($\hat{\beta}_{\mathrm{rev}} \approx -2.24$).
\end{tablenotes}
\end{table}

\begin{figure}[htbp]
\centering
\includegraphics[width=0.84\textwidth]{welfare_theta.pdf}
\caption{Welfare loss as a function of the
  diagnostic-overreaction parameter $\theta$ in the
  calibrated New Keynesian benchmark.}
\label{fig:welfare_theta}
\end{figure}

The welfare schedule is convex in $\theta$: partial
reductions in overreaction yield first-order gains when
the economy operates in a high-$\theta$ region, even if
communication policy cannot eliminate behavioral
distortions entirely.

Two caveats apply. First, the mapping from a reduced-form
revision--error coefficient to a structural $\theta$
parameter requires assumptions about the data-generating
process that this paper does not formally test. The
calibration is therefore a plausibility check, not a
structural identification result. Second, if policy
communication changes household belief formation, the
reduced-form $\theta$ can shift endogenously (Lucas
critique). The reported numbers are local counterfactuals
around the estimated regime, not invariant policy
constants.

